% REV02: pruebas reunidas en 02_dualidad_t.tex y 03_pantallas_coordinadas.tex.
\subsection{Incidencia excepcional y coordinación de las pantallas}
\label{exc:pantallas-dualidad}

La dirección seleccionada por \(K\) y el retículo de Leech proceden
de representaciones distintas de la incidencia construida. La primera
determina la reducción ortogonal \(12\to11\to10\); la segunda
interviene en la reducción isotrópica \(26\to24\). La coordinación
conserva el espacio previo, sus relaciones y la información transmitida
por cada proyección.

Las secciones \ref{iv:dualidad} y \ref{iv:pantallas} desarrollan
conjuntamente esa coordinación. La primera reúne la involución de
coordenadas enteras, la conjugación unitaria con sus dominios y el
cambio de sección residual con memoria del cociente. La segunda
construye el morfismo \(\mathcal R_M^{\mathrm{alg}}\) de
\eqref{iv:eq:rm-alg} y demuestra el isomorfismo de las presentaciones
cocientadas en el teorema \ref{iv:thm:pantallas-isomorfas}.
El grafo conserva la coordenada en \(V_{11}\) junto con su proyección
en \(V_{10}\), mientras el cociente isotrópico recupera el retículo
de Leech. Así se preservan los respectivos núcleos y codominios.

La elipse de acción se construye antes de aplicar la dualidad y fija
su radio adimensional junto con el radio recíproco
(sección \ref{iv:ellipse:section}). El sector compacto enlaza entonces
la incidencia excepcional con una forma cuadrática determinada. Los
campos, la acción y los operadores de la realización física posterior
se presentan en su dominio propio, conservando este antecedente
algebraico y sus demostraciones internas.
